\section{INNER JOIN, LEFT JOIN, RIGHT JOIN}

\textbf{JOIN} adalah operasi untuk menggabungkan baris dari dua atau lebih tabel berdasarkan kolom yang berhubungan (biasanya foreign key dan primary key). \code{INNER JOIN} mengembalikan hanya baris yang memiliki pasangan di \textbf{kedua} tabel. \code{LEFT JOIN} mengembalikan semua baris dari tabel kiri dan baris yang cocok dari tabel kanan; jika tidak ada pasangan, kolom dari tabel kanan bernilai NULL. \code{RIGHT JOIN} sebaliknya: semua baris dari tabel kanan \cite{mysqlDocs}.

\begin{figure}[ht]
\centering
\begin{tikzpicture}[
  circle/.style={draw, minimum size=2cm, font=\small},
  node distance=0.5cm
]
  % INNER JOIN
  \begin{scope}[xshift=0cm]
    \node[circle, fill=blue!10] (a1) {A};
    \node[circle, fill=green!10, right=-0.7cm of a1] (b1) {B};
    \begin{scope}
      \clip (a1.center) circle (1cm);
      \fill[orange!30] (b1.center) circle (1cm);
    \end{scope}
    \draw (a1.center) circle (1cm);
    \draw (b1.center) circle (1cm);
    \node[below=0.8cm of a1, xshift=0.35cm, font=\footnotesize\bfseries] {INNER JOIN};
  \end{scope}

  % LEFT JOIN
  \begin{scope}[xshift=4.5cm]
    \node[circle, fill=orange!30] (a2) {A};
    \node[circle, fill=green!10, right=-0.7cm of a2] (b2) {B};
    \begin{scope}
      \clip (a2.center) circle (1cm);
      \fill[orange!30] (b2.center) circle (1cm);
    \end{scope}
    \draw (a2.center) circle (1cm);
    \draw (b2.center) circle (1cm);
    \node[below=0.8cm of a2, xshift=0.35cm, font=\footnotesize\bfseries] {LEFT JOIN};
  \end{scope}

  % RIGHT JOIN
  \begin{scope}[xshift=9cm]
    \node[circle, fill=blue!10] (a3) {A};
    \node[circle, fill=orange!30, right=-0.7cm of a3] (b3) {B};
    \begin{scope}
      \clip (b3.center) circle (1cm);
      \fill[orange!30] (a3.center) circle (1cm);
    \end{scope}
    \draw (a3.center) circle (1cm);
    \draw (b3.center) circle (1cm);
    \node[below=0.8cm of a3, xshift=0.35cm, font=\footnotesize\bfseries] {RIGHT JOIN};
  \end{scope}
\end{tikzpicture}
\caption{Visualisasi Venn diagram untuk jenis-jenis JOIN}
\end{figure}

\begin{table}[ht]
\centering
\begin{tabular}{|P{3cm}|P{9.5cm}|}
\hline
\textbf{Jenis JOIN} & \textbf{Baris yang Dikembalikan} \\
\hline
\code{INNER JOIN} & Hanya baris yang ada pasangannya di kedua tabel \\
\hline
\code{LEFT JOIN} & Semua baris tabel kiri + pasangan dari kanan (NULL jika tidak ada) \\
\hline
\code{RIGHT JOIN} & Semua baris tabel kanan + pasangan dari kiri (NULL jika tidak ada) \\
\hline
\code{CROSS JOIN} & Semua kombinasi baris (cartesian product) \\
\hline
\end{tabular}
\caption{Perbandingan jenis-jenis JOIN}
\end{table}

\begin{webcode}[language=SQL, caption={Contoh penggunaan JOIN}]
-- INNER JOIN: produk beserta nama kategorinya
SELECT p.nama, p.harga, k.nama AS kategori
FROM produk p
INNER JOIN kategori k ON p.kategori_id = k.id;

-- LEFT JOIN: semua produk, termasuk yang belum punya kategori
SELECT p.nama, p.harga, IFNULL(k.nama, 'Tanpa Kategori') AS kategori
FROM produk p
LEFT JOIN kategori k ON p.kategori_id = k.id;

-- Multi-table JOIN
SELECT p.nama AS produk, k.nama AS kategori, u.nama AS pembeli
FROM pesanan_detail pd
INNER JOIN produk p ON pd.produk_id = p.id
INNER JOIN kategori k ON p.kategori_id = k.id
INNER JOIN pesanan ps ON pd.pesanan_id = ps.id
INNER JOIN pengguna u ON ps.user_id = u.id;
\end{webcode}

